\begin{textsample}{POS Dim 2 – ms – Score 22.00 – ms/ms\_em\_1.txt}  \label{ex:f2_pos_075}
APRESENTAÇÃO O \textbf{Governo} do Estado de Mato Grosso do Sul, por meio da Secretaria de Estado de Educação ( SED/MS ), apresenta o \textbf{Currículo} de Referência de Mato Grosso do Sul – Etapa do Ensino Médio.
Este documento surge como parte do processo de implementação desta etapa da educação básica, decorrente da Lei n. 13.415, de 16 de \textbf{fevereiro} de 2017, que altera a Lei 9.394/1996 - \textbf{Diretrizes} e Bases da Educação Nacional ( LDB ).
Partindo do pressuposto de que o \textbf{Currículo} de Referência contempla as expectativas locais para a formação dos estudantes, sua construção ocorreu de forma colaborativa com a sociedade sul-mato-grossense, com vistas ao desenvolvimento das aprendizagens essenciais, enriquecidas pelo contexto histórico, econômico, ambiental, cultural e do mundo do trabalho e da prática social vivenciada no Estado.
A Lei n. 13.415/2017, que alterou a Lei n. 9.394, de 20 de dezembro de 1996, incorporou novas características ao Ensino Médio de todo o país.
Por conseguinte, o Conselho Nacional de Educação ( CNE ), por meio da Resolução CNE/CEB n. 03, de 21 de \textbf{novembro} de 2018, \textbf{atualizou} as \textbf{Diretrizes} Curriculares Nacionais para o Ensino Médio.
Nesse ínterim, a Resolução CNE/CP n. 04, de 17 de dezembro de 2018, \textbf{homologou} a Base Nacional Comum Curricular – etapa do Ensino Médio ( BNCC-EM ), o que possibilitou ao Estado de Mato Grosso do Sul iniciar a elaboração deste documento.
Com o apoio do Ministério da Educação ( MEC ), em consonância com a \textbf{Portaria} n. 1.371, de 16 de \textbf{julho} de 2019, \textbf{publicada} no Diário Oficial da União, em 17 de \textbf{julho} de 2020, foi constituída uma \textbf{equipe} representativa[^1 ] da SED/MS e da UNDIME ( União dos \textbf{Dirigentes} \textbf{Municipais} de Educação ) para atuar no processo de implementação desse \textbf{Currículo}, e é composta pelos seguintes integrantes : Coordenador Estadual, Coordenador de Etapa, Articulador entre Etapas, Articulador de Itinerário Formativo Propedêutico, Articulador de Itinerário da Formação \textbf{Técnica} e Profissional, Coordenadores de Áreas e \textbf{redatores} formadores.
A entrega deste documento configura uma conquista para todas as redes públicas e \textbf{instituições} privadas pertencentes ao Sistema Estadual de Ensino de Mato Grosso do Sul, pois remonta à importância das parcerias institucionais no âmbito do regime de colaboração.
Desse modo, o \textbf{Currículo} de Referência do Estado de Mato Grosso do Sul – Etapa do Ensino Médio constitui o documento normativo para a compreensão, adequação e \textbf{qualificação} do Projeto Político Pedagógico ( PPP ) das unidades escolares e a organização do trabalho didático dos professores com vistas à formação integral dos estudantes. [ ^1 ] : Além das funções já identificadas, a \textbf{equipe} representativa também é composta por integrantes da Educação Infantil e do Ensino Fundamental, bem como de articulador do Conselho Estadual de Educação/MS.

% matched lemmas: atualizar, currículo, diretriz, dirigente, equipe, fevereiro, governo, homologar, instituição, julho, municipal, novembro, portaria, publicar, qualificação, redator, técnico
\end{textsample}
